\import{test/phase7/composition-generalized-base.tex}

\begin{abbreviation}\label{calculation_parameterized_equality}
    $a\calceq{n}b$ iff $(a,n)=(b,n)$.
\end{abbreviation}

\begin{composition}\label{calculation_parameterized_composition}
    Suppose $x\calceq{n}y\calceq{n}z$.
    Then $x\calceq{n}z$.
\end{composition}
\begin{proof}
    Follows.
\end{proof}

\begin{composition}\label{calculation_negative_membership}
    Suppose $x\notin B\supseteq A$.
    Then $x\notin A$.
\end{composition}
\begin{proof}
    Follows by \obv{subseteq}.
\end{proof}

\begin{abbreviation}\label{calculation_guarded_equality}
    $a\calcguard{A}b$ iff $(a,A)=(b,A)$ and $A=A$.
\end{abbreviation}

\begin{composition}\label{calculation_several_conditions}
    Suppose $x\in A$.
    Suppose $y\in A$ and $z\in A$.
    Suppose $x\in A$.
    Suppose $x\calcguard{A}y\calcguard{A}z$.
    Then $x\calcguard{A}z$.
\end{composition}
\begin{proof}
    Follows.
\end{proof}

\begin{abbreviation}\label{calculation_reversed_membership}
    $a\calcreverse b$ iff $b\in a$.
\end{abbreviation}

\begin{abbreviation}\label{calculation_ignored_endpoint}
    $a\calcignore b$ iff $a=a$.
\end{abbreviation}

\begin{abbreviation}\label{calculation_erased_parameter}
    $a\calcerase{n}b$ iff $a\subseteq b$.
\end{abbreviation}

\begin{definition}\label{calculation_typed_relation}
    Let $x\in A$.
    $x\calctyped{A}y$ iff $x=y$.
\end{definition}
